\section{Itô 积分与随机微分方程}
\label{sec:06-Random-Processes/06-Diffusion-and-SDE/03}

前两节已经大咧咧地写着 \(dX_t=\mu\,dt+\sigma\,dW_t\)，而\hyperref[sec:06-Random-Processes/06-Diffusion-and-SDE/01]{``从随机游走到布朗运动''}一节早已警告：布朗路径处处不可微，总变差无界。\(dW\) 不是普通微分，把这个记号当真写进方程，等于在普通微积分的语法里塞了一个不存在的东西。这一节来还这笔账：先定义 \(\int\sigma\,dW\) 到底是什么意思，随机微分方程（SDE）才不是一串好看的符号。

\subsection{一个试图定义积分的努力，以及它为什么失败}

面对 \(\int_0^T f\,dW\)，最自然的思路是借用 Riemann--Stieltjes 的构造：把区间剖细，取点，求和，取极限。对普通积分 \(\int f\,dg\)，只要 \(f\) 连续、\(g\) 有界变差，这个极限存在，而且跟取样点 \(\xi_k\) 选在区间左端、右端还是中间无关。

布朗运动呢？\hyperref[sec:06-Random-Processes/06-Diffusion-and-SDE/01]{``从随机游走到布朗运动''}一节已经算清楚了一件关键事实：位移的量级是 \(\sqrt{\Delta t}\) 而不是 \(\Delta t\)。这个 \(\sqrt{\Delta t}\) 尺度的后果是冷酷的：\(\sum|\Delta W_k|\) 量级约为 \(\sqrt{\Delta t}\cdot(1/\Delta t)=1/\sqrt{\Delta t}\)，剖分越细，和越大——以概率一，布朗路径在任意区间上的总变差无界。

Riemann--Stieltjes 积分对无界变差的积分器不存在——不是技术困难，是根本不存在。按普通微积分的标准，\(\int \cdot\,dW\) 没有意义。需要另一条路。

可是，同一个布朗路径的增量平方之和却收敛：\(\sum(\Delta W_k)^2\to T\)。位移的绝对值之和发散，位移的平方之和却老实收敛到区间长度——这就是二次变差 \([W]_t=t\)。两条性质来自同一个 \(\sqrt{\Delta t}\) 尺度，一条堵死了逐路径积分的路，另一条给出了唯一的可控量。

\subsection{换一个框架：不在每条路径上定义积分}

既然路径太粗糙，就不能要求积分值由单条路径上的极限确定。Itô 的策略是在均方意义下定义积分——只要求积分的二阶矩行为良好，不要求每条路径上的 Riemann 和都收敛。这个转向是整台机器的支点。

构造分三步，本书给出核心思想和关键估计，完备化（取极限时的闭包论证）依赖于测度论，对此声明存在、指明条件即可。

第一步，对阶梯状的简单被积过程定义积分。设
\[
H_t=\sum_k H_k\,\ind_{[t_k,t_{k+1})}(t),
\]
其中 \(H_k\) 在 \(t_k\) 时刻就已经确定——它只能依赖 \(t_k\) 及之前的信息，不能偷看区间内部或未来的布朗运动。这称为适应性（adaptivity），也叫非预见性（non-anticipating）。对这个过程，直接定义
\[
\int_0^T H\,dW=\sum_k H_k\big(W_{t_{k+1}}-W_{t_k}\big).
\]
这里的关键选择是权重 \(H_k\) 在区间\textbf{左端点}取值，乘上它之后发生的增量。先下注，后开奖——这一步的选择会在后面引发连锁后果。

第二步，这个积分的二阶矩有一个简洁的公式。计算平方期望时，交叉项的条件期望为零（因为 \(H_k\) 已确定，\(\Delta W_k\) 独立且均值为零），剩下
\[
\E\big[H_k^2(\Delta W_k)^2\big]=\E[H_k^2]\,\Delta t_k.
\]
对所有区间求和得 Itô 等距：
\[
\E\left[\left(\int_0^T H\,dW\right)^2\right]
=\E\left[\int_0^T H_t^2\,dt\right].
\]
左边是积分的二阶矩，右边是被积函数平方的普通 Lebesgue 积分的时间平均。这个等式说，映射 \(H\mapsto\int H\,dW\) 保持 \(L^2\) 距离。等式右端出现 \(dt\) 而不是 \(dW\) 的平方，这恰恰是二次变差 \(\E[(\Delta W)^2]=\Delta t\) 的功劳——若没有这个收敛事实，右边的 \(dt\) 就没有来由。

第三步，保持距离意味着极限可以传递。任何能被阶梯过程在 \(L^2\) 意义下逼近的适应过程，其 Itô 积分定义为阶梯过程积分的均方极限。这个极限存在且唯一，不依赖逼近序列的选择——Itô 等距保证了这两点。这就是 Itô 积分。

因为是均方意义下的极限，Itô 积分天然携带 martingale 性质：在适当的可积条件下，积分过程关于上限是 martingale，且
\[
\E\left[\int_0^T H\,dW\right]=0.
\]
公平性来自左端点取值：权重确定的时刻，它所乘的增量尚未发生。这与\hyperref[sec:06-Random-Processes/04-Applications/02]{``Martingale：公平游戏与停时''}中``没有可预测漂移''的结构一脉相承。

\subsection{代价：链式法则多出一项}

普通积分不在乎取样点——左端点、右端点、中点，只要剖分够细，极限都一样。Itô 积分不行。取样点不同，极限也不同，因为权重与区间内的位移相关，相关量级是 \(\sqrt{\Delta t}\cdot\sqrt{\Delta t}=\Delta t\)——剖分再细也不消失。

Itô 选择了左端点，代价是普通微积分的链式法则不再成立。对 \(f(X_t)\) 做 Taylor 展开：
\[
df=f'(X)\,dX+\frac12 f''(X)\,(dX)^2+\cdots
\]
在普通微积分中，\((dX)^2\) 是 \((dt)^2\) 量级，可以丢掉。在 Itô 的世界里，\(dX\) 含有 \(\sigma\,dW\)，而 \((dW)^2\) 是 \(dt\) 量级——不是高阶无穷小，不能丢。这条乘法规则必须记住：
\[
(dt)^2=0,\qquad dt\,dW=0,\qquad (dW)^2=dt.
\]
前两条来自 \(dt\) 的相对尺度过小，最后一条则是二次变差 \([W]_t=t\) 的微分写法——它是这整个修正项的来源。

代入 \(dX_t=\mu\,dt+\sigma\,dW_t\)，展开到二阶：
\[
\begin{aligned}
df(X_t)&=f'(X_t)(\mu\,dt+\sigma\,dW_t)+\frac12 f''(X_t)(\mu\,dt+\sigma\,dW_t)^2\\
&=f'(X_t)\mu\,dt+f'(X_t)\sigma\,dW_t+\frac12 f''(X_t)\sigma^2\,dt.
\end{aligned}
\]
整理得 Itô 公式：
\[
\boxed{\;df(X_t)=\left(f'(X_t)\,\mu+\frac12f''(X_t)\,\sigma^2\right)dt
+f'(X_t)\,\sigma\,dW_t\;}
\]
修正项 \(\frac12 f''\sigma^2\) 是``\(\sqrt{\Delta t}\) 世界''的入场费。只要被变换的函数有曲率，噪声的二阶波动就会渗入一阶漂移。第一次见到这个公式时最容易犯的错，就是把 \((dW)^2\) 当成 \((dt)^2\) 丢掉——本节的乘法规则正是为了拦住这一步。

\subsection{动手算一个：几何布朗运动}

最简单的随机价格模型假设相对涨落恒定：
\[
dS_t=\mu S_t\,dt+\sigma S_t\,dW_t.
\]
要把这个式子积出显式解，需要对 \(S_t\) 做一个变换，让噪声从乘法变成加法。试试取对数 \(f(S)=\log S\)。\(f'(S)=1/S\)，\(f''(S)=-1/S^2\)。代入 Itô 公式：
\[
\begin{aligned}
d\log S_t&=\left(\frac{1}{S_t}\cdot\mu S_t+\frac12\cdot\left(-\frac{1}{S_t^2}\right)\cdot\sigma^2 S_t^2\right)dt+\frac{1}{S_t}\cdot\sigma S_t\,dW_t\\
&=\left(\mu-\frac{\sigma^2}{2}\right)dt+\sigma\,dW_t.
\end{aligned}
\]
漂移系数从 \(\mu\) 变成了 \(\mu-\sigma^2/2\)，多了一个负的修正。积分得
\[
S_t=S_0\exp\!\left(\left(\mu-\frac{\sigma^2}{2}\right)t+\sigma W_t\right).
\]
如果用普通链式法则，会猜指数里是 \(\mu t+\sigma W_t\)，漏掉了 \(- \sigma^2t/2\)。

这个修正有实在的含义。计算期望：由 Gaussian 矩母函数 \(\E[e^{\sigma W_t}]=e^{\sigma^2t/2}\)，得 \(\E[S_t]=S_0 e^{\mu t}\)——期望增长率确实是 \(\mu\)。但每条路径的对数增长率只有 \(\mu-\sigma^2/2\)，比 \(\mu\) 少了一块。差距就是波动侵蚀（volatility drag）：典型的路径被波动拖着往下走，但少数暴涨的路径把期望拉了回去。Jensen 不等式在指数函数上作用的结果，正是 Itô 修正项的金融面孔。

\subsection{再算一个：Ornstein--Uhlenbeck 过程}

上一节用 Fokker--Planck 的密度方法研究过均值回复模型
\[
dX_t=-\theta X_t\,dt+\sigma\,dW_t,\qquad\theta>0.
\]
现在用路径方法直接解它。Itô 公式有一个直截了当的推广，当 \(f\) 显含时间时：
\[
df(X_t,t)=\left(\frac{\partial f}{\partial t}
+\mu\,\frac{\partial f}{\partial x}
+\frac12\sigma^2\frac{\partial^2 f}{\partial x^2}\right)dt
+\sigma\frac{\partial f}{\partial x}\,dW.
\]
对 \(\mu=-\theta X\)、常数 \(\sigma\)，取 \(f(x,t)=e^{\theta t}x\)（相当于 ODE 解法中的积分因子）。\(f\) 对 \(x\) 二阶导为零，修正项消失：
\[
d\big(e^{\theta t}X_t\big)=e^{\theta t}dX_t+\theta e^{\theta t}X_t\,dt
=\sigma e^{\theta t}\,dW_t.
\]
积分得显式解
\[
X_t=X_0e^{-\theta t}+\sigma\int_0^t e^{-\theta(t-s)}\,dW_s.
\]
被积函数 \(e^{-\theta(t-s)}\) 是确定性的，不依赖 \(W\)，所以 Itô 积分结果是 Gaussian 随机变量。均值是 \(X_0e^{-\theta t}\)——初值的影响按指数衰减。方差由 Itô 等距给出：
\[
\sigma^2\int_0^t e^{-2\theta(t-s)}\,ds
=\frac{\sigma^2}{2\theta}\big(1-e^{-2\theta t}\big)
\;\longrightarrow\;\frac{\sigma^2}{2\theta}.
\]
极限分布 \(\mathcal N(0,\sigma^2/(2\theta))\) 与上一节用 Fokker--Planck 方程求出的平稳密度一致。路径视角与密度视角在同一个极限处汇合——两个完全不同的数学工具给了同一个答案。初值被均值回复的力量按 \(e^{-\theta t}\) 的速度遗忘，噪声的累计效应被稳定在有限方差内，方差不再增长。

\subsection{SDE 是积分方程的缩写}

记号 \(dX=\mu\,dt+\sigma\,dW\) 虽然方便，但严格含义只能是积分方程
\[
X_t=X_0+\int_0^t\mu(X_s)\,ds+\int_0^t\sigma(X_s)\,dW_s.
\]
第一项是普通 Lebesgue 积分（或 Riemann 积分，漂移通常足够光滑），第二项是 Itô 积分。``微分形式''是缩写；检验一个解，永远要代回积分形式。

解的存在唯一性有标准的充分条件：若 \(\mu\) 和 \(\sigma\) 关于 \(x\) 满足 Lipschitz 连续，且线性增长
\[
\abs{\mu(x)}+\abs{\sigma(x)}\le C(1+\abs{x}),
\]
则强解存在且唯一（声明，不提供完整证明）。线性增长条件的作用是排除有限时间爆炸——确定性方程 \(dx=x^2\,dt\) 已经展示超线性漂移可以在有限时间内把解送到无穷。噪声有时会压制爆炸（扩散把质量带离危险区域），有时也可能诱发爆炸（扩散把质量推进危险区域），不能离开具体的 \(\mu,\sigma\) 一概而论。状态依赖的 \(\sigma\) 碰到零、带有跳跃项、有反射边界等情况都需要各自的理论，已超出本书范围。

\subsection{Stratonovich 的另一套选择}

区间中点取值给出另一套积分，称为 Stratonovich 积分，记作 \(\sigma\circ dW\)。它恢复了普通形式的链式法则，在几何和物理的坐标变换中更自然——物理文献常把噪声理解为光滑噪声的极限，极限路径刚好对应 Stratonovich。代价是积分一般不再是 martingale。

两种解释之间相差一个漂移修正。一维时，
\[
\sigma\circ dW=\sigma\,dW+\frac12\,\sigma\frac{\partial\sigma}{\partial x}\,dt.
\]
其中 \(\circ\) 标记 Stratonovich。若噪声强度依赖状态（即 \(\sigma\) 不是常数），修正项出现；若 \(\sigma\) 是常数，两者一致。

金融与本书采用 Itô，因为``先下注后开奖''的左端点结构保持了 martingale 性质，定价与对冲中的期望计算因此简化。两种体系内部各自自洽，混用才会出错。本单元此后只用 Itô。

解析可解的 SDE 屈指可数——几何布朗运动和 OU 过程正是屈指可数的那几个。下一篇转向数值模拟：既然大多数 SDE 写不出显式解，我们至少要能可靠地算出近似路径。数值模拟还会引出一个反直觉的方向——时间反演，沿着它我们将瞥见扩散生成模型的门径。
